\chapter{Biểu diễn tri thức và Trạng thái}

\section{Định nghĩa không gian trạng thái}

Trạng thái toàn cục $S$ của một ván cờ tướng tại một thời điểm được định nghĩa chính thức bởi bộ bốn thành phần:
\[ S = \langle B, p, c_h, h \rangle \]

Trong đó:
\begin{itemize}
    \item $B$ (Board) đại diện cho phân bố các quân cờ trên bàn lưới kích thước $10 \times 9 = 90$ ô. Mỗi ô cờ tại vị trí $sq \in [0, 89]$ có thể trống hoặc chứa một trong 14 loại quân cờ thực tế (Tướng, Sĩ, Tượng, Xe, Pháo, Mã, Tốt cho mỗi bên Đỏ và Đen).
    \item $p$ (Player/Side) đại diện cho bên đi tiếp theo, $p \in \{\text{Đỏ}, \text{Đen}\}$.
    \item $c_h$ (Count Half-move) là số nửa nước đi (ply) trôi qua kể từ nước đi không có bắt quân cuối cùng. Đây là bộ đếm phục vụ cho việc thực thi luật hòa cờ 60 nước.
    \item $h$ (History) là danh sách lịch sử lưu trữ toàn bộ trạng thái di chuyển trước đó cùng các thông tin bắt quân, hỗ trợ hoàn tác nước đi và phát hiện các trường hợp lặp lại trạng thái (repetition check).
\end{itemize}

\section{Kỹ thuật biểu diễn trạng thái trong Lingine}

Để hiện thực hóa mô hình toán học trên với hiệu năng cao nhất, Lingine định nghĩa cấu trúc dữ liệu trung tâm \texttt{Position} và cấu trúc thông tin phụ trợ \texttt{StateInfo} trong tệp \texttt{src/core/position/mod.rs}:

\begin{lstlisting}[language=Rust, caption=Cấu trúc Position trong Lingine]
#[derive(Clone)]
pub struct Position {
    // Flat 90-square array mapping square index (0 to 89) to the Piece occupying it.
    board: [Option<Piece>; Square::COUNT],
    // Precomputed bitboards showing piece placements grouped by PieceType.
    bitboard_by_type: [Bitboard; PieceType::COUNT],
    // Precomputed bitboards showing piece placements grouped by Side.
    bitboard_by_color: [Bitboard; Side::COUNT],

    // Active count of each piece category on the board.
    piece_count: [u8; Piece::COUNT],
    // Current state of position
    state: StateInfo,
    // Stack tracking previous move parameter histories for undoing moves.
    history: Vec<StateInfo>,
    // Total moves played in the game so far
    game_ply: u16,
    // Palace coordinates of both sides' Generals (Kings) for faster check detection.
    king_squares: [Square; Side::COUNT],
}
\end{lstlisting}

Để tiết kiệm thời gian tính toán lại thế cờ trong quá trình tìm kiếm sâu, các dữ liệu động phát sinh sau mỗi nước đi được đóng gói riêng vào cấu trúc \texttt{StateInfo}:

\begin{lstlisting}[language=Rust, caption=Cấu trúc StateInfo trong Lingine]
#[derive(Clone, Copy, Debug)]
pub struct StateInfo {
    // The last move played to reach this state.
    pub last_move: Option<Move>,
    // The piece captured during this move, or None if it was quiet.
    pub captured_piece: Option<Piece>,
    // The prior Zobrist position hash value before the move occurred.
    pub zobrist: u64,
    // Halfmove clock / 60-rule counter
    pub sixtymove_clock: u16,
    // Whether the side to move was in check in this position state.
    pub in_check: bool,
    // Precalculated incremental mid- and end-game score (from Red's perspective)
    pub score: PackedScore,
    // Precalculated incremental game phase
    pub phase: u8,
    // Checker pieces checking the King of the side to move.
    pub checkers: Bitboard,
    // Pinned pieces (blockers) for both sides' kings.
    pub blockers_for_king: [Bitboard; Side::COUNT],
    // The slider pieces pinning other pieces for both sides' kings.
    pub pinners: [Bitboard; Side::COUNT],
    // Check squares for each piece type of the side to move
    pub check_squares: [Bitboard; PieceType::COUNT],
    // Whether we need a full check validation.
    pub need_full_check: bool,
    // Number of plies played since the last null move.
    pub plies_since_null: u16,
}
\end{lstlisting}

\subsection{Biểu diễn bàn cờ bằng Bitboard 128-bit}

Do bàn cờ tướng có kích thước $10 \times 9 = 90$ ô, cấu trúc Bitboard chuẩn 64-bit của cờ vua ($8 \times 8 = 64$ ô) là hoàn toàn không đủ. Lingine đã phát triển một giải pháp Bitboard tự chọn dựa trên kiểu số nguyên không dấu 128-bit (\texttt{u128}), trong đó 90 bit đầu tiên được gán tương ứng cho 90 ô cờ từ vị trí $0$ (a1) tới $89$ (i10):
\[ \text{Bitboard} = \sum_{sq=0}^{89} b_{sq} \cdot 2^{sq} \quad (\text{với } b_{sq} \in \{0, 1\}) \]

Cấu trúc \texttt{Bitboard} được đóng gói trong tệp \texttt{src/core/bitboard.rs} như sau:

\begin{lstlisting}[language=Rust, caption=Cấu trúc Bitboard trong Lingine]
#[derive(Copy, Clone, Default, PartialEq, Eq, Debug)]
pub struct Bitboard(u128);

impl Bitboard {
    #[inline]
    pub const fn raw(self) -> u128 {
        self.0
    }
    
    // Loai bo bit bat (1) thap nhat va tra ve chi so cua no
    #[inline]
    pub fn pop_lsb(&mut self) -> Option<Square> {
        if self.0 == 0 {
            None
        } else {
            let lsb = self.0.trailing_zeros() as u8;
            self.0 &= self.0 - 1;
            Some(unsafe { std::mem::transmute(lsb) })
        }
    }
    
    // Dem so bit bat (1)
    #[inline]
    pub const fn count_ones(self) -> u32 {
        self.0.count_ones()
    }
}
\end{lstlisting}

Việc sử dụng \texttt{u128} cho phép CPU thực hiện các phép tính logic bit song song trên toàn bộ 90 ô cờ chỉ bằng một chỉ thị máy duy nhất. Nhờ cơ chế \texttt{pop\_lsb} sử dụng lệnh assembly \texttt{TZCNT} của vi xử lý x86-64 thông qua hàm \texttt{trailing\_zeros()}, việc duyệt qua tất cả các quân cờ đang hoạt động đạt tốc độ tối đa.

\subsection{Biểu diễn các thành phần lượt đi và lịch sử}

\begin{itemize}
    \item \textbf{Lượt đi ($p$):} Được biểu diễn bằng kiểu \texttt{Side} gồm hai giá trị \texttt{Side::Red} và \texttt{Side::Black}. Các phương thức đối lập \texttt{Side::opposite()} giúp đổi lượt nhanh bằng phép toán số học đơn giản.
    \item \textbf{Bộ đếm luật hòa cờ ($c_h$):} Được lưu trong trường \texttt{sixtymove\_clock: u16}. Mỗi nước đi quiet (không ăn quân) sẽ tăng bộ đếm lên 1. Nếu có quân bị bắt, bộ đếm lập tức reset về 0. Khi giá trị đạt 120 (tương đương 60 nước đi đầy đủ của cả hai bên), động cơ sẽ nhận diện thế hòa.
    \item \textbf{Lịch sử ($h$):} Lưu trữ trong stack \texttt{history: Vec<StateInfo>}. Mỗi khi thực hiện nước đi thông qua \texttt{do\_move()}, trạng thái \texttt{StateInfo} cũ được đẩy vào \texttt{history}. Khi hoàn tác qua \texttt{undo\_move()}, trạng thái chỉ cần pop từ stack ra để khôi phục nhanh chóng trong thời gian $O(1)$ mà không cần tính toán lại từ đầu.
\end{itemize}

\section{Nén dữ liệu nước đi 16-bit}

Trong quá trình duyệt cây tìm kiếm, hàng triệu nước đi được tạo ra và truyền qua lại giữa các hàm. Nếu sử dụng các cấu trúc dữ liệu cồng kềnh, chi phí cấp phát bộ nhớ và truyền tham số sẽ làm chậm đáng kể tốc độ của động cơ. Lingine giải quyết triệt để vấn đề này bằng cách nén nước đi thành một số nguyên không dấu 16-bit (\texttt{u16}) duy nhất:

\begin{table}[h]
    \centering
    \begin{tabular}{ccl}
        \toprule
        \textbf{Các bit} & \textbf{Khoảng giá trị} & \textbf{Ý nghĩa} \\
        \midrule
        0 -- 6 & 0 .. 89 & Chỉ số ô đích (To Square) -- Cần 7 bit ($2^7 = 128$) \\
        7 -- 13 & 0 .. 89 & Chỉ số ô xuất phát (From Square) -- Cần 7 bit \\
        14 -- 15 & 0 .. 3 & Cờ phân loại nước đi đặc biệt (Flags) \\
        \bottomrule
    \end{tabular}
    \caption{Cấu trúc mã hóa nước đi 16-bit trong Lingine}
\end{table}

\begin{lstlisting}[language=Rust, caption=Cấu trúc Move trong Lingine]
#[derive(Copy, Clone, PartialEq, Eq, Debug)]
pub struct Move(u16);

impl Move {
    #[inline]
    pub const fn new(from: Square, to: Square) -> Self {
        Self(((from as u16) << 7) | (to as u16))
    }
    
    #[inline]
    pub const fn from(self) -> Square {
        unsafe { std::mem::transmute((self.0 >> 7) & 0x7F) }
    }
    
    #[inline]
    pub const fn to(self) -> Square {
        unsafe { std::mem::transmute(self.0 & 0x7F) }
    }
}
\end{lstlisting}

\subsection{Ví dụ minh họa}
Giả sử ta cần mã hóa nước đi di chuyển quân Xe từ ô xuất phát $h2$ (chỉ số 16, nhị phân \texttt{0010000}) đến ô đích $e2$ (chỉ số 13, nhị phân \texttt{0001101}):
\begin{itemize}
    \item Chỉ số ô đi (From): $16 \ll 7 = 2048$ (nhị phân \texttt{00001000 00000000})
    \item Chỉ số ô đến (To): $13$ (nhị phân \texttt{00000000 00001101})
    \item Giá trị nước đi mã hóa: $\text{Move} = 2048 \lor 13 = 2061$ (thập lục phân \texttt{0x080D})
\end{itemize}

Nhờ kích thước cực nhỏ (2 bytes), dữ liệu nước đi có thể nằm gọn trong các thanh ghi CPU và truyền trực tiếp bằng trị số, tránh tối đa việc truy xuất bộ nhớ RAM và tận dụng tối ưu bộ nhớ đệm L1 Cache.

\section{Định danh trạng thái bằng Zobrist Hashing}

Zobrist Hashing là kỹ thuật sinh mã định danh 64-bit gần như duy nhất cho mỗi trạng thái bàn cờ. Bằng cách gán cho mỗi cặp (loại quân cờ, ô đứng) và lượt đi một số ngẫu nhiên 64-bit, mã băm Zobrist giúp kiểm tra lặp thế trận cực nhanh và làm khóa tra bảng chuyển vị (Transposition Table).

\subsection{Công thức băm}

Mã băm của trạng thái thế cờ $S_0$ được tính bằng:
\[ \text{Hash}(S_0) = Z_{side} \oplus \left( \bigoplus_{sq=0}^{89} Z_{piece(sq), sq} \right) \]

Trong đó:
\begin{itemize}
    \item $Z_{side}$ là giá trị ngẫu nhiên 64-bit đại diện cho bên được quyền đi.
    \item $Z_{piece(sq), sq}$ là giá trị ngẫu nhiên 64-bit tương ứng với quân cờ đang ngự trị tại ô $sq$.
\end{itemize}

\subsection{Khởi tạo bảng băm ngẫu nhiên tĩnh tại thời điểm biên dịch}

Để đảm bảo tính nhất quán và tốc độ tối đa, toàn bộ bảng mã ngẫu nhiên Zobrist được tính toán trực tiếp tại thời điểm biên dịch (Compile-time) thông qua cơ chế \texttt{const fn} của Rust. Nhóm phát triển đã thiết lập gieo hạt (Seed) cho bộ sinh số ngẫu nhiên LCG (Linear Congruential Generator) dựa trên việc thực hiện phép toán XOR mã số sinh viên của cả 3 thành viên trong nhóm:
\[ \text{SEED} = 202416124 \oplus 202400076 \oplus 2416167 = 203875323 \]

Mã nguồn khởi tạo tại \texttt{src/core/position/zobrist\_table.rs}:

\begin{lstlisting}[language=Rust, caption=Khởi tạo Zobrist Table tĩnh tại thời điểm biên dịch]
pub struct ZobristTable {
    pub pieces: [[u64; Square::COUNT]; Piece::COUNT],
    pub side: u64,
}

pub const ZOBRIST: ZobristTable = init_zobrist();

const fn init_zobrist() -> ZobristTable {
    let mut seed = 202416124 ^ 202400076 ^ 2416167; // Seed tu MSSV
    let mut lcg = || {
        seed = seed.wrapping_mul(6364136223846793005).wrapping_add(1442695040888963407);
        seed
    };
    
    let mut table = ZobristTable {
        pieces: [[0; Square::COUNT]; Piece::COUNT],
        side: 0,
    };
    // Sinh khoa ngau nhien cho tung quan co tai moi o co
    ...
    table.side = lcg();
    table
}
\end{lstlisting}

\subsection{Cập nhật vi phân trạng thái $O(1)$}

Khi di chuyển quân cờ, mã băm Zobrist không cần phải quét duyệt lại toàn bộ bàn cờ. Nhờ tính chất của phép XOR ($A \oplus A = 0$ và $0 \oplus B = B$), ta có thể cập nhật vi phân mã băm chỉ bằng các phép XOR tại các vị trí thay đổi:

\begin{lstlisting}[language=Rust, caption=Cập nhật vi phân Zobrist trong do\_move]
// Loai bo quan co tai o di (from) khoi ma bam
hash ^= ZOBRIST.pieces[piece][from];

// Them quan co vao o den (to) trong ma bam
hash ^= ZOBRIST.pieces[piece][to];

// Neu la nuoc di an quan, loai bo quan co bi bat ra khoi ma bam
if let Some(captured) = captured_piece {
    hash ^= ZOBRIST.pieces[captured][to];
}

// Thay doi luot di
hash ^= ZOBRIST.side;
\end{lstlisting}

Cơ chế cập nhật vi phân này biến thao tác tính toán mã băm từ độ phức tạp $O(N)$ (với $N=90$ ô cờ) về độ phức tạp $O(1)$ chỉ với 3 đến 4 phép XOR cơ bản, tiết kiệm tối đa tài nguyên CPU.

\section{Bộ phân xử luật cờ tướng và phát hiện lặp trạng thái (Rule Judge)}

Luật chơi cờ tướng Việt Nam và quốc tế có điểm khác biệt rất lớn so với cờ vua trong việc xử lý các trạng thái lặp lại thế trận (Repetitions). Trong cờ vua, nếu một thế trận lặp lại 3 lần (Threefold Repetition), ván đấu lập tức được xử hòa. Tuy nhiên, trong cờ tướng, luật chơi nghiêm cấm các hành vi mang tính chất cưỡng đoạt hoặc quấy rối vô hạn:
\begin{itemize}
    \item \textbf{Chiếu dai (Perpetual Check):} Một bên liên tục chiếu tướng đối phương mà không thay đổi phương án tấn công. Bên chiếu dai bắt buộc phải thay đổi nước đi, nếu không sẽ bị xử thua.
    \item \textbf{Đuổi dai (Perpetual Chase):} Một bên liên tục sử dụng các quân tấn công (Xe, Mã, Pháo, Tốt đã qua sông) để uy hiếp, bắt một quân cờ của đối phương một cách lặp đi lặp lại. Bên đuổi dai cũng bị cấm và sẽ bị xử thua nếu không đổi nước đi.
    \item \textbf{Lặp nước bình thường:} Các tình huống hai bên luân phiên đổi nước đi tự vệ, hoặc quấy rối lẫn nhau (không vi phạm chiếu dai hay đuổi dai) sẽ được xử hòa.
\end{itemize}

Lingine tích hợp bộ trọng tài luật \texttt{RuleJudge} trong tệp \texttt{src\slash core\slash position\slash rule\_judge.rs}. Để phát hiện lặp lại thế trận, hệ thống lưu trữ toàn bộ lịch sử các mã băm Zobrist trong danh sách \texttt{history}. Khi tìm kiếm, động cơ sẽ duyệt ngược lịch sử để tìm các mã băm trùng lặp. Khi phát hiện trùng lặp thế trận ở ply thứ $i$, bộ trọng tài sẽ kích hoạt giải thuật kiểm tra đuổi bắt (\texttt{chased}) để xác định hành vi của bên đi:

\begin{lstlisting}[language=Rust, caption=Giải thuật xác định hành vi đuổi quân trong Lingine]
fn chased(&mut self, mover: Side, id_board: &[u8; Square::COUNT]) -> u16 {
    let mut chase = 0u16;
    let opponent = mover.opposite();
    let occupied = self.bitboard_occupied();

    // 1. Quet cac quan muc tieu co the bi duoi cua doi phuong (tru Tuong)
    let mut targets_mask = self.bitboard_by_type[PieceType::Rook]
        | self.bitboard_by_type[PieceType::Cannon]
        | self.bitboard_by_type[PieceType::Knight]
        | self.bitboard_by_type[PieceType::Advisor]
        | self.bitboard_by_type[PieceType::Bishop];

    // Them cac quan Tot da qua song cua doi thu
    let opp_pawns = self.bitboard_by_type[PieceType::Pawn] & self.bitboard_by_color[opponent];
    targets_mask |= opp_pawns & Bitboard::side(mover);

    let targets = targets_mask & self.bitboard_by_color[opponent];
    ...
}
\end{lstlisting}

Giải thuật kiểm tra đuổi bắt sẽ thực hiện di chuyển thử nghiệm ảo trên bàn cờ. Nếu quân tấn công có giá trị thấp hơn hoặc bằng quân bị tấn công, hoặc quân bị tấn công đang đứng ở vị trí không được bảo vệ (unprotected), nước đi đó sẽ được phân loại là hành vi ``Đuổi'' (\texttt{Chase}). Nếu một bên thực hiện liên tiếp các nước đi chiếu hoặc đuổi lặp lại trong chu kỳ thế trận, bộ trọng tài sẽ trả về điểm số phạt thua đối với bên vi phạm:
\begin{itemize}
    \item Nếu bên đi phạm luật chiếu dai hoặc đuổi dai: Trả về điểm số bị chiếu bí sau $d$ nước đi ($\text{mated\_in}(ply)$).
    \item Nếu đối thủ phạm luật: Trả về điểm số chiếu bí đối thủ ($\text{mate\_in}(ply)$).
    \item Nếu lặp nước hòa hợp lệ: Trả về điểm số hòa cờ ($\text{score::ZERO} = 0$).
\end{itemize}

\section{Thiết kế bảng chuyển vị Transposition Table (TT)}

Bảng chuyển vị (Transposition Table -- TT) là bộ nhớ đệm cực kỳ quan trọng giúp ghi nhớ kết quả đánh giá và nước đi tốt nhất tại các nút đã duyệt. Do cây tìm kiếm cờ tướng có rất nhiều đường dẫn trùng lặp (hoán vị nước đi, ví dụ đi Xe trước rồi đi Mã sau so với đi Mã trước Xe sau), TT giúp cắt tỉa các nhánh cờ trùng lặp này ngay tại gốc.

\subsection{Cấu trúc dữ liệu tối ưu hóa bộ nhớ đệm}

Để tối ưu hóa việc truy xuất bộ nhớ, cấu trúc lưu trữ của bảng băm \texttt{Storage} được thiết kế với kích thước cực kỳ nhỏ gọn, vừa vặn đúng 16 byte (128 bit). Điều này giúp các entry nằm thẳng hàng trong các đường truyền cache line của CPU, giảm thiểu tối đa hiện tượng cache miss:

\begin{lstlisting}[language=Rust, caption=Cấu trúc Storage 16-byte của bảng băm TT]
#[derive(Clone, Debug)]
struct Storage {
    pub key: Key,         // 8 bytes: Khoa Zobrist Hash
    pub score: Score,     // 4 bytes: Diem so danh gia the co
    pub best_move: Option<Move>, // 2 bytes: Nuoc di tot nhat
    pub flags: Flags,     // 1 bytes: Loai bound (Exact/Alpha/Beta) + Age
    pub depth: i8,        // 1 bytes: Do sau tim kiem da thuc hien
}
\end{lstlisting}

Nhờ sử dụng kiểu dữ liệu nguyên tử \texttt{NonZeroU8} cho \texttt{flags}, trình biên dịch Rust thực hiện tối ưu hóa ngách (niche optimization) để nén toàn bộ thực thể \texttt{Option<Storage>} vào đúng 16 byte mà không cần thêm byte padding nào.

\subsection{Chính sách thay thế ưu tiên Độ sâu và Tuổi (Replacement Strategy)}

Khi xảy ra va chạm chỉ số bảng băm (hai thế trận khác nhau ánh xạ vào cùng một chỉ số trong bảng), Lingine áp dụng chính sách thay thế kết hợp **Độ sâu** (Depth-preferred) và **Tuổi** (Age-preferred):
\begin{itemize}
    \item \textbf{Độ sâu:} Các kết quả tìm kiếm ở độ sâu lớn hơn chứa nhiều thông tin chiến thuật hơn, do đó được ưu tiên giữ lại.
    \item \textbf{Tuổi (Age):} Mỗi lần thực hiện tìm kiếm mới ở nước đi tiếp theo, chỉ số tuổi của bảng băm (\texttt{age}) tăng lên. Các entry quá cũ (từ các nước đi trước đó) sẽ giảm giá trị tham chiếu và dễ bị ghi đè hơn.
\end{itemize}

Công thức kiểm tra điều kiện ghi đè entry trong \texttt{src\slash search\slash transposition\_table.rs} được cài đặt như sau:
\begin{lstlisting}[language=Rust]
let is_better = value.depth >= existing.depth - self.relative_age(existing.flags.age()) as i8;
if is_better {
    self.table[index] = Some(Storage { ... });
}
\end{lstlisting}

\subsection{Chuyển đổi điểm số Mate phụ thuộc độ sâu (Mate Distance Pruning)}

Một vấn đề kỹ thuật phức tạp khi sử dụng bảng băm là việc lưu trữ điểm số chiếu sát (Mate score). Điểm số chiếu sát trả về từ bộ tìm kiếm luôn phụ thuộc vào số nước đi từ gốc tới vị trí đó ($ply$). Nếu ta lưu trực tiếp điểm số này vào TT, khi một nhánh tìm kiếm khác ở độ sâu khác truy vấn trúng entry này, điểm số chiếu sát sẽ bị lệch nước đi.

Lingine giải quyết triệt để bằng cách chuẩn hóa điểm số trước khi lưu và khôi phục khi truy vấn:
\begin{itemize}
    \item \textbf{Khi lưu (Store):} Chuyển đổi điểm số từ phụ thuộc ply sang độc lập ply:
    \[ \text{Score}_{\text{TT}} = \begin{cases} 
      \text{Score} + ply & (\text{nếu thắng chiếu sát}) \\
      \text{Score} - ply & (\text{nếu thua chiếu sát}) 
   \end{cases} \]
    \item \textbf{Khi truy vấn (Probe):} Khôi phục lại điểm số tương thích với ply hiện tại:
    \[ \text{Score}_{\text{Search}} = \begin{cases} 
      \text{Score}_{\text{TT}} - ply & (\text{nếu thắng chiếu sát}) \\
      \text{Score}_{\text{TT}} + ply & (\text{nếu thua chiếu sát}) 
   \end{cases} \]
\end{itemize}
